% year: תשע"ט
% semester: ב
% moed: ב
% number: 1ג
% points: 5
% categories: func-analysis
% source: exams/מבחנים/תשעט/חודיסמן מועד ב תשעט.pdf

\begin{question}
עבור אילו ערכים של $a$ יש לקו $y=x^2+a\sin x$ נקודות פיתול?
\end{question}

\begin{hint}
בדקו מתי $y''$ מחליפה סימן.
\end{hint}

\begin{solution}
$y'=2x+a\cos x$, $y''=2-a\sin x$. נקודת פיתול היא נקודה שבה $y''$ מחליפה סימן.
\begin{itemize}
  \item אם $|a|<2$: $|a\sin x|\le|a|<2$ ולכן $y''>0$ לכל $x$ – אין נקודות פיתול.
  \item אם $a=2$: $y''=2(1-\sin x)\ge0$, ומתאפסת רק בנקודות מבודדות; הסימן לא מתחלף – אין פיתול. באופן דומה עבור $a=-2$: $y''=2(1+\sin x)\ge0$.
  \item אם $|a|>2$: המשוואה $\sin x=\frac2a$ עם $0<\left|\frac2a\right|<1$ יש לה פתרונות $x_0$, ובהם $\cos x_0\ne0$. לכן $(y'')'(x_0)=-a\cos x_0\ne0$, כלומר $y''$ עוברת דרך $0$ בצורה מונוטונית ממש ומחליפה סימן ב-$x_0$ – נקודת פיתול.
\end{itemize}
\textbf{תשובה:} $|a|>2$.
\end{solution}
